\section{The whole original geometric kernel module maps into the entire rounded pushforward}
\label{sec:rounded-canonical-kernel-module-map}

\begin{definition}[The same whole pushed canonical differential section]
\label{sc0-rounded-kernel-map:original-pushed-section}
\lean{EnvelopeProperty.Char0.RoundedCanonicalModule.pushedDifferentialSection,EnvelopeProperty.Char0.RoundedCanonicalModule.pushedDifferentialSection_app,EnvelopeProperty.Char0.RoundedCanonicalModule.pushedDifferentialSection_quotientProjection}
\leanok
Adjoint the SAME original canonical differential Hom-section by
the actual pullback--pushforward adjunction and original pullback
unit comparison. Its value on EVERY original local function is
the SAME original differential applied to the original pullback
of that function. Projection by the SAME pushed rounded quotient
is exactly the pre-existing original pushed quotient section.
\end{definition}
\begin{proof}
Use the true native unit adjunction on the SAME whole section.
The proved original unit-section evaluation identity retains
every original section value. Actual adjunction naturality,
the original quotient-section definition, genuine functor
composition and associativity give the entire original projection
identity. No section-value or matching formula is assumed.
\end{proof}

\begin{theorem}[The original kernel module lifts into the entire pushed rounded module]
\label{sc0-rounded-kernel-map:whole-original-lift}
\lean{EnvelopeProperty.Char0.RoundedCanonicalModule.kernelToPushforwardRounded,EnvelopeProperty.Char0.RoundedCanonicalModule.kernelToPushforwardRounded_inclusion}
\leanok
Let the original source and target be smooth of the same relative
dimension over the original field, let $g$ be quasi-compact, and
let the original ideal have its actual SNC divisor condition.
The SAME whole geometric rounded canonical kernel ideal MODULE
has an actual whole map to the ENTIRE original pushed rounded
relative canonical module. Composing with the SAME pushed rounded
inclusion gives exactly the SAME original ideal inclusion followed
by the original pushed canonical differential section.
No pushforward-quasicoherence, vanishing, affine, flatness or
factorization certificate is supplied.
\end{theorem}
\begin{proof}
Actual smoothness derives monicity of the SAME original rounded
inclusion. Its actual native cokernel sequence is left exact.
The original pushforward is a genuine right adjoint and additive,
so it preserves the relevant actual kernels, monomorphisms and
zero morphisms. Thus the SAME original pushed rounded sequence
is left exact, without needing whole pushforward quasi-coherence.

The actual quotient is quasicoherent, derived from the original
smooth/SNC tensor geometry. The preceding whole pushed-kernel
ideal MODULE theorem shows its actual original inclusion is
annihilated by the SAME original pushed quotient section.
The preceding whole projection identity therefore proves that
the original ideal inclusion followed by the SAME pushed canonical
differential lies in the actual pushed rounded kernel. The native
left-exact kernel lift gives the asserted entire original module
map. Its genuine lift identity is exactly the stated full
inclusion comparison, retaining every original map.
\uses{sc0-rounded-kernel-map:original-pushed-section,
      sc0-pushed-kernel-module:whole-native-kernel}
\end{proof}

This constructs the genuine whole comparison map. Bijectivity
and its identification with the SAME full all-DVR multiplier
module remain the next consumers; no vanishing is inferred.
